The formal artifact verifies the algebraic construction underlying
Theorem~\ref{thm:unified}. Its shared definitions import Mathlib; the
construction proof is a separate module. The challenge imports only
those definitions, and the submitted solution imports the construction
proof. The sources and pinned project are in the \artifactlink{}.

\lstdefinestyle{leancode}{basicstyle=\ttfamily\footnotesize,
  keywordstyle=\bfseries,backgroundcolor=\color{achieve!5},
  frame=single,rulecolor=\color{achieve!35},framerule=.4pt,
  xleftmargin=4pt,xrightmargin=4pt,framexleftmargin=4pt,
  breaklines=true,columns=fullflexible,keepspaces=true,
  showstringspaces=false,escapeinside={(*@}{@*)},
  morekeywords={theorem,by,exact,let,have,rcases,with,obtain,if,then,else}}

\subsection{Exact goal and mathematical interpretation}\label{app:lean-theorem}
Use the field $\F_q$ and enumeration $u_1,\ldots,u_q$ from
Appendix~\ref{app:theorem}. Lean denotes this field by $\mathbb F$ and
its cardinality by \texttt{Fintype.card} $\mathbb F$, which equals $q$.
The natural number \texttt{tail} is $t=n-2q>0$; hence $n=2q+t$.
The disjunction below specifies the same $q\ge3$ and odd-length binary
cases as Theorem~\ref{thm:unified}. Within namespace \texttt{ICLR2027EAQECC.Comparator}, the
submitted declaration is:

% Verbatim excerpt generated by scripts/make_lean_excerpt.py.
\begin{lstlisting}[style=leancode]
theorem theorem1_exact_goal
    {(*@$\mathbb{F}$@*) : Type u} [Field (*@$\mathbb{F}$@*)] [Fintype (*@$\mathbb{F}$@*)] [DecidableEq (*@$\mathbb{F}$@*)]
    {tail : (*@$\mathbb{N}$@*)} [NeZero tail]
    (hcase : 3 (*@$\le$@*) Fintype.card (*@$\mathbb{F}$@*) (*@$\lor$@*)
      (Fintype.card (*@$\mathbb{F}$@*) = 2 (*@$\land$@*) Odd (2 * Fintype.card (*@$\mathbb{F}$@*) + tail))) :
    (*@$\exists$@*) z(*@${}_2$@*) lam : (*@$\mathbb{F}$@*), lam (*@$\ne$@*) 0 (*@$\land$@*) 1 + z(*@${}_2$@*) (*@$\ne$@*) 0 (*@$\land$@*)
      Nonempty (NormalizerSideCertificate (*@$\mathbb{F}$@*) tail z(*@${}_2$@*) lam) := by
  exact theorem1_exact hcase
\end{lstlisting}


Lean represents the radical coefficients by the function space
\[
 \mathcal T=\{f:\F_q\to\F_q\}.
\]
A function $f$ assigns one scalar to each field element; $f$ need not be linear. Addition and scalar multiplication are pointwise,
$(f+g)(u)=f(u)+g(u)$ and $(a f)(u)=a f(u)$ for $a\in\F_q$.
The enumeration $u_1,\ldots,u_q$ identifies $f$ with the vector
$\vct\tau=(f(u_1),\ldots,f(u_q))\in\F_q^q$.
This evaluation map is a linear isomorphism, so $\dim_{\F_q}\mathcal T=q$.
Thus Lean's \texttt{ClosedCoeff} type is
$\mathcal C=\F_q^2\oplus\mathcal T\cong\F_q^{q+2}$.
Write $\vct\eta=(\alpha,\beta,f)$ and $\tau_i=f(u_i)$, matching the
coefficients of $\vct a,\vct b,\vct\rho_i$ in Appendix~\ref{app:theorem}.
Relabeling the paired coordinates by $i=1,\ldots,q$, the linear map
$\Phi=\texttt{closedWord}$ is
\[
\begin{aligned}
 \Phi(\vct{\eta})_{(i,0)}&=(\alpha+\beta u_i,\,\beta+\tau_i),\\
 \Phi(\vct{\eta})_{(i,1)}&=(\alpha+\beta u_i,\,\beta z_2-\tau_i),\\
 \Phi(\vct{\eta})_{\mathrm{tail},j}&=(\alpha,\,\beta\lambda_j),
 \qquad 0\le j<t,
\end{aligned}
\]
where $\lambda_0=\lambda$ and $\lambda_j=1$ for $j>0$.
Thus $\Phi(\vct{\eta})=\alpha\vct a+\beta\vct b+
\sum_{i=1}^q\tau_i\vct\rho_i$ in the formal coordinate order.
Let $L=\operatorname{im}\Phi$, $\mathcal C_0=
\{(0,0,f):f\in\mathcal T\}$, and
$\delta=\langle\Phi(1,0,0_{\mathcal T}),\Phi(0,1,0_{\mathcal T})\rangle_s$,
where $0_{\mathcal T}$ is the zero function.
The certificate constructed by \texttt{theorem1\_exact} proves:
\begingroup
\renewcommand{\labelenumi}{(\arabic{enumi})}
\begin{enumerate}\setlength{\itemsep}{2pt}
\item \texttt{encodeInjective} and \texttt{deltaNonzero}:
$\Phi$ is injective and $\delta\ne0$.
\item \texttt{gramFormula} and \texttt{radicalIff}: for all
$\vct\eta=(\alpha,\beta,f)$ and
$\vct\eta'=(\alpha',\beta',g)$,
\[
 \langle\Phi(\vct\eta),\Phi(\vct\eta')\rangle_s
 = (\alpha\beta'-\beta\alpha')\delta,
 \qquad \Phi(\vct\eta)\in\operatorname{rad}(L)
 \iff \vct\eta\in\mathcal C_0.
\]
\item \texttt{normalizerFinrank} and \texttt{radicalFinrank}:
$\dim L=q+2$ and $\dim\mathcal C_0=q$.
\item \texttt{distanceLower} and \texttt{distanceAttained}:
\[
 \forall\vct\eta\notin\mathcal C_0,\quad
 \wt(\Phi(\vct\eta))\ge n-1,
 \qquad
 \exists\vct\eta\notin\mathcal C_0,\quad
 \wt(\Phi(\vct\eta))=n-1.
\]
\end{enumerate}
\endgroup
Injectivity identifies $\Phi(\mathcal C_0)$ with the $q$-dimensional
radical of $L$. Taking $S=L^\perp$ gives $\dim S=2n-q-2$.
The standard stabilizer correspondence
(Appendix~\ref{app:symplectic}) therefore yields
$c=(\dim S-q)/2=n-q-1$, $k=n-\dim S+c=1$ and $d=n-1$.
These are the code parameters recorded by the certificate; the
proof-bearing fields above certify their underlying symplectic construction.

For the coordinate convention, Lean's $(u_i,\mathrm{false})$ and
$(u_i,\mathrm{true})$ are our
$(i,0)$ and $(i,1)$. Its tail puts $\lambda$ first; Appendix~\ref{app:theorem}
puts it last. This permutation preserves the symplectic product and weight.
The proof chooses $z_2=0$ in characteristic two and $z_2=1$ otherwise.

Figure~\ref{fig:lean-dependencies} groups the certificate's mathematical
dependencies. The dashed arrow is the standard stabilizer interpretation
used in the paper; it is separate from the Lean-checked algebraic fields.
% Grouped mathematical dependencies, checked against Theorem1.lean.
\begin{figure}[!htbp]
\centering
\begin{tikzpicture}[font=\footnotesize,
  box/.style={draw=black!30,rounded corners=2pt,fill=black!3,
    align=center,inner sep=4pt,minimum height=8mm},
  arr/.style={-{Latex[length=1.7mm]},draw=black!55,line width=.6pt}]
\node[box,text width=12.1cm] (map) at (0,0)
  {Coordinate map $\Phi$ and admissible parameters
   $\lambda\ne0$, $1+z_2\ne0$, $\delta\ne0$};
\node[box,text width=3.65cm,minimum height=12mm] (dim) at (-4.15,-1.6)
  {Injectivity\\$\dim L=q+2$\\$\dim\mathcal C_0=q$};
\node[box,text width=3.65cm,minimum height=12mm] (gram) at (0,-1.6)
  {Gram formula\\Radical identification\\$\operatorname{rad}(L)=\Phi(\mathcal C_0)$};
\node[box,text width=3.65cm,minimum height=12mm] (dist) at (4.15,-1.6)
  {Weight lower bound\\$\wt\ge n-1$\\Equality attained};
\node[box,fill=achieve!8,draw=achieve!55,text width=5.3cm] (cert) at (-2.6,-3.35)
  {Lean algebraic certificate\\\texttt{NormalizerSideCertificate}};
\node[box,fill=white,text width=5.1cm] (code) at (3.5,-3.35)
  {Standard stabilizer correspondence\\$[[n,1,n-1;n-q-1]]_q$};
\draw[arr] (map.south -| dim.north)--(dim.north);
\draw[arr] (map.south)--(gram.north);
\draw[arr] (map.south -| dist.north)--(dist.north);
\draw[arr] (dim.south)--(dim.south |- cert.north);
\draw[arr] (gram.south)--(cert.north);
\draw[arr] (dist.south)--([xshift=1.9cm]cert.north);
\draw[arr,dashed] (cert.east)--(code.west);
\end{tikzpicture}
\caption{Dependencies of the algebraic certificate.}
\label{fig:lean-dependencies}
\end{figure}

\FloatBarrier

\subsection{Solution and formal verification}\label{app:lean-replay}
The solution above invokes \texttt{theorem1\_exact}. Its proof body is
reproduced verbatim below. For $q\ge3$, it chooses $\lambda$ outside two
forbidden values. In the binary branch, odd length permits $\lambda=1$.
Both branches invoke \texttt{theorem1\_closed\_form}, which assembles the
certificate from the injectivity, Gram, radical, dimension and distance
lemmas in the same source module.

% Verbatim excerpt generated by scripts/make_lean_excerpt.py.
\begin{lstlisting}[style=leancode]
  let z(*@${}_2$@*) : (*@$\mathbb{F}$@*) := if (2 : (*@$\mathbb{F}$@*)) = 0 then 0 else 1
  have hz(*@${}_2$@*) : 1 + z(*@${}_2$@*) (*@$\ne$@*) 0 := by
    dsimp [z(*@${}_2$@*)]
    split_ifs with htwo
    (*@$\cdot$@*) simp
    (*@$\cdot$@*) simpa [one_add_one_eq_two] using htwo
  rcases hcase with hlarge | (*@$\langle$@*)hcard, hodd(*@$\rangle$@*)
  (*@$\cdot$@*) obtain (*@$\langle$@*)lam, hlam, h(*@$\delta$@*)base(*@$\rangle$@*) := exists_nonzero_add_ne_zero hlarge
        (((2 * Fintype.card (*@$\mathbb{F}$@*) + tail - 1 : (*@$\mathbb{N}$@*)) : (*@$\mathbb{F}$@*)))
    have h(*@$\delta$@*) : closedDelta (tail := tail) z(*@${}_2$@*) lam (*@$\ne$@*) 0 := by
      rw [closedDelta_eq_cast_length_sub_one_add]
      exact h(*@$\delta$@*)base
    exact (*@$\langle$@*)z(*@${}_2$@*), lam, hlam, hz(*@${}_2$@*), theorem1_closed_form z(*@${}_2$@*) lam hlam hz(*@${}_2$@*) h(*@$\delta$@*)(*@$\rangle$@*)
  (*@$\cdot$@*) let lam : (*@$\mathbb{F}$@*) := 1
    have hlam : lam (*@$\ne$@*) 0 := one_ne_zero
    have h(*@$\delta$@*) : closedDelta (tail := tail) z(*@${}_2$@*) lam (*@$\ne$@*) 0 := by
      rw [closedDelta_eq_cast_length_sub_one_add]
      exact binary_odd_delta_ne_zero hcard hodd
    exact (*@$\langle$@*)z(*@${}_2$@*), lam, hlam, hz(*@${}_2$@*), theorem1_closed_form z(*@${}_2$@*) lam hlam hz(*@${}_2$@*) h(*@$\delta$@*)(*@$\rangle$@*)
\end{lstlisting}


The proof and submitted solution use only \texttt{propext},
\texttt{Quot.sound} and \texttt{Classical.choice}. Comparator confirms
that the solution proves the specified goal under this axiom policy,
and Lean's kernel accepts the proof~\citep{LeanComparator}.
The \artifactlink{} provides the pinned project and reproducible
verification records.
